\documentclass{article}

% 修改为 final 选项：去掉 preprint 页脚，并显示作者姓名
\usepackage[final]{neurips_2026}

\usepackage[utf8]{inputenc} % allow utf-8 input
\usepackage[T1]{fontenc}    % use 8-bit T1 fonts
\usepackage{hyperref}       % hyperlinks
\usepackage{url}            % simple URL typesetting
\usepackage{booktabs}       % professional-quality tables
\usepackage{amsfonts}       % blackboard math symbols
\usepackage{nicefrac}       % compact symbols for 1/2, etc.
\usepackage{microtype}      % microtypography
\usepackage{xcolor}         % colors

\title{Project Proposal: Building and Extending GPT-2}

\author{%
  Han Huang \\
  \texttt{hanhuang@kth.se} \\
  \And
  Jingze Guo\\
  \texttt{jingze@kth.se} \\
  \AND
  Kai Zhao\\
  \texttt{kaizha@kth.se} \\
  \And
  Boyan Zhang \\
  \texttt{boyanz@kth.se} \\
}

\begin{document}

\maketitle

\section{Project Type}
% 3) Whether you will complete a custom, default or default + extension project.
Default Project 3b (with A/B extensions).

\section{Problem Description and Proposed Approach}
% 4) A brief description of the problem that you will work on and how you will try to solve it. Reference to at least one paper.

\paragraph{Problem Description} 
While large language models (LLMs) based on the Transformer architecture demonstrate exceptional zero-shot and few-shot capabilities, adapting these massive pre-trained models to specific downstream tasks remains computationally expensive and memory-intensive. The primary objective of this project is to gain a rigorous, low-level understanding of the Transformer architecture by building the core components of GPT-2 and exploring how to effectively fine-tune it under constrained computational resources.

\paragraph{Proposed Approach} 
We will solve this by implementing a bare-bones version of OpenAI’s GPT-2 architecture and the Adam optimizer from scratch. To bypass the prohibitively expensive pre-training phase, we will initialize our custom architecture using pre-trained weights from Hugging Face.

Our baseline approach involves extracting representations from the frozen model to train a linear classifier, which we will then compare against full model fine-tuning. We will evaluate the model's performance on sentiment classification using the Stanford Sentiment Treebank (SST) and CFIMDB datasets, and adapt it for text generation.

Furthermore, to address the high memory footprint of full fine-tuning and as our primary extension for a higher grade, we will investigate parameter-efficient fine-tuning (PEFT). Specifically, we propose implementing Low-Rank Adaptation (LoRA) to update only a small fraction of the model's weights, comparing its final performance and computational efficiency against the standard full fine-tuning approach.

\paragraph{Inspirational Literature} 
This project’s foundation is inspired by the original GPT-2 paper, ``Language Models are Unsupervised Multitask Learners'' \cite{radford2019gpt2}. Our proposed extension methodology is directly guided by the LoRA paper \cite{hu2021lora}.


\section{Data}
% 5) The data that you will use for training, validation and testing.
Stanford Sentiment Treebank (SST) and CFIMDB for sentiment analysis; Shakespeare's Sonnets for text generation.

\section{Software Packages}
% 6) The deep learning software package(s) you will use.
PyTorch

\section{Implementation Details}
% 7) How much of the software implementation your group will write and to what extent will rely on open-source implementations.
[Specify what you will code from scratch. e.g., We will write the missing core Transformer blocks and Adam optimizer, while relying on Hugging Face for the pre-trained weights.]

\section{Baselines}
% 8) The baselines you intend to use.
[Specify baselines. e.g., Linear classifier trained on extracted representations; un-finetuned GPT-2 performance.]

\section{Evaluation Metrics}
% 9) How will you evaluate your results?
[Specify metrics. e.g., Accuracy for classification tasks; Perplexity and BLEU scores for text generation.]

\section{Target Grade}
% 10) The grade your project group is aiming for in the range A-E.
A

\section{Milestones}
% 11) Which achieved milestones during your project should be reached as you work towards your grade target.

\subsection{E/C Grade Milestones}
% 去掉了 \item 后面的中括号
\begin{itemize}
    \item Add milestone 1, e.g., Write and debug the 5 missing code blocks.
    \item Add milestone 2, e.g., Fine-tune the pre-trained model for sentiment analysis on SST and CFIMDB.
    \item Add milestone 3, e.g., Fine-tune for text generation.
\end{itemize}

\subsection{C/D Grade Milestones}
% 去掉了 \item 后面的中括号
\begin{itemize}
    \item Add milestone 1, e.g., Perform thorough hyper-parameter search.
    \item Add milestone 2, e.g., Compare full fine-tuning versus training a linear classifier on extracted representations.
\end{itemize}

\subsection{A/B Grade Milestones}
% 去掉了 \item 后面的中括号
\begin{itemize}
    \item Add milestone 1, e.g., Implement an extension such as parameter-efficient fine-tuning (LoRA), accelerating attention, or model quantization.
    \item Add milestone 2, e.g., Conduct ablation studies on the chosen extension.
\end{itemize}

\section{Skills to Acquire}
% 12) Specify for each group member the skills/knowledge w.r.t. deep learning "theory" and practice they aim to acquire.
\begin{itemize}
    \item \textbf{Han Huang:} [Skill to acquire]
    \item \textbf{Jingze Guo:} [Skill to acquire]
    \item \textbf{Kai Zhao:} [Skill to acquire]
    \item \textbf{Boyan Zhang:} [Skill to acquire]
\end{itemize}

\section{Workload Breakdown (3hp per person)}
% 13) As the project corresponds to 3hp per person, please specify how the proposed project will require this amount of work per person.


% --- References Section ---
\bibliographystyle{plain}
\begin{thebibliography}{10}

\bibitem{radford2019gpt2}
Radford, A., Wu, J., Child, R., Luan, D., Amodei, D., \& Sutskever, I.
\newblock Language models are unsupervised multitask learners.
\newblock \emph{OpenAI blog}, 1(8):9, 2019.

\bibitem{hu2021lora}
Hu, E. J., Shen, Y., Wallis, P., Allen-Zhu, Z., Li, Y., Wang, S., Wang, L., \& Chen, W.
\newblock LoRA: Low-rank adaptation of large language models.
\newblock \emph{arXiv preprint arXiv:2106.09685}, 2021.

\end{thebibliography}

\end{document}